\documentclass[a4paper,11pt]{article}
\usepackage[a4paper,margin=2cm]{geometry}
\usepackage[T1]{fontenc}
\usepackage[utf8]{inputenc}
\usepackage[ngerman,english]{babel}
\usepackage{lmodern,microtype,xcolor,parskip,enumitem}
\usepackage[most]{tcolorbox}
\usepackage{titlesec}
\usepackage[hidelinks]{hyperref}
\definecolor{HeaderBlueDark}{RGB}{70,110,170}
\definecolor{RowBlueLight}{RGB}{235,243,255}
\definecolor{RowBlueDark}{RGB}{220,233,250}
\definecolor{SoftGray}{RGB}{90,90,90}
\setlength{\parindent}{0pt}
\setlist[itemize]{leftmargin=1.4em,itemsep=0.35em,topsep=0.35em}
\titleformat{\section}[block]{\Large\bfseries\color{HeaderBlueDark}}{}{0pt}{}
\newtcolorbox{exbox}{enhanced,breakable,colback=RowBlueLight,colframe=RowBlueDark,boxrule=0.4pt,arc=1.5mm,left=1.2mm,right=1.2mm,top=1mm,bottom=1mm,title={Beispiele \textnormal{(Examples)}},fonttitle=\bfseries,colbacktitle=RowBlueDark,coltitle=black}
\NewDocumentEnvironment{grammarentry}{mm+b}{\begin{tcolorbox}[enhanced,breakable,colback=white,colframe=HeaderBlueDark,title={#1\hfill\normalfont\footnotesize #2},colbacktitle=HeaderBlueDark,coltitle=white,fonttitle=\bfseries]#3\end{tcolorbox}}{}
\newcommand{\GermanText}[1]{\textbf{Deutsch: }#1\par}
\newcommand{\EnglishText}[1]{{\small\color{SoftGray}\textbf{English: }#1\par}}
\newcommand{\ExampleItem}[2]{\item \textbf{#1}\\{\small\color{SoftGray}#2}}
\title{Die Modalpartikel \glqq denn\grqq\\{\large\textnormal{The modal particle denn}}}
\date{}
\begin{document}
\maketitle
\tableofcontents
\section{Grundbedeutung \textnormal{(Basic meaning)}}
\begin{grammarentry}{denn}{Modalpartikel / Abtönungspartikel}
\GermanText{In Fragen macht \textbf{denn} die Frage gesprächsbezogener und drückt je nach Kontext Neugier, Überraschung, Ungeduld oder Vorwurf aus. Es ist kein Fragewort und verlangt keinen Kasus.}
\EnglishText{In questions, \textbf{denn} conveys conversational interest and may express curiosity, surprise, impatience or reproach. It is not an interrogative word and governs no case.}
\end{grammarentry}
\section{Satzart und Stellung \textnormal{(Sentence type and position)}}
\begin{grammarentry}{Mittelfeld}{Fragepartikel}
\GermanText{\textbf{denn} ist als Modalpartikel meist unbetont und steht im Mittelfeld von W-Fragen und Ja/Nein-Fragen. Verkürzte Fragen wie \glqq Was denn?\grqq\ sind elliptisch.}
\EnglishText{As a modal particle, \textbf{denn} is usually unstressed and appears in the middle field of wh-questions and yes/no questions. Short questions like \glqq Was denn?\grqq\ are elliptical.}
\end{grammarentry}
\begin{exbox}\begin{itemize}
\ExampleItem{Was machst du denn?}{So, what are you doing?}
\ExampleItem{Wo warst du denn?}{Where have you been, then?}
\ExampleItem{Hast du denn Zeit?}{Do you actually have time?}
\ExampleItem{Warum kommst du denn nicht?}{Why aren't you coming, then?}
\ExampleItem{Was ist denn bloß passiert?}{What on earth happened?}
\end{itemize}\end{exbox}
\section{denn nicht \textnormal{(denn with not)}}
\begin{grammarentry}{Zwei Wortarten}{Modalpartikel + Negationspartikel}
\GermanText{In \glqq Warum kommst du denn nicht?\grqq\ ist \textbf{denn} eine Modalpartikel und \textbf{nicht} eine Negationspartikel. Es handelt sich \textbf{nicht} um eine einzige grammatische Einheit. Die Wörter dürfen getrennt stehen.}
\EnglishText{In \glqq Warum kommst du denn nicht?\grqq, \textbf{denn} is a modal particle and \textbf{nicht} negates the sentence. They \textbf{do not} form one grammatical unit and need not be adjacent.}
\end{grammarentry}
\begin{exbox}\begin{itemize}
\ExampleItem{Warum hast du denn gestern nicht angerufen?}{Why didn't you call yesterday, then?}
\ExampleItem{Warum kommst du nicht?}{Why aren't you coming? (without the nuance of denn)}
\end{itemize}\end{exbox}
\section{Andere grammatische Funktion \textnormal{(Other grammatical role)}}
\begin{grammarentry}{denn als Konjunktion}{weil / because}
\GermanText{\textbf{denn} kann auch eine nebenordnende Kausalkonjunktion sein. Dann verbindet es Hauptsätze und bedeutet ungefähr \glqq weil\grqq. Als Modalpartikel verbindet es dagegen keine zwei Sätze.}
\EnglishText{\textbf{denn} can also be a coordinating causal conjunction meaning approximately \glqq because\grqq. In that use it joins independent clauses, unlike the modal particle.}
\end{grammarentry}
\begin{exbox}\begin{itemize}
\ExampleItem{Ich bleibe zu Hause, denn ich bin krank.}{I am staying home because I am ill.}
\ExampleItem{Warum bleibst du denn zu Hause?}{So why are you staying at home?}
\end{itemize}\end{exbox}
\section{Quellen \textnormal{(Sources)}}
\begin{itemize}
\item IDS Mannheim: \url{https://grammis.ids-mannheim.de/terminologie/2}
\item Duden: \url{https://www.duden.de/sprachwissen/sprachratgeber/Modalpartikeln}
\end{itemize}
\end{document}
